\chapter{Engine Implementation and Configuration}
\label{app:config}

This appendix explains how the two matching engines are built: how an
order is represented, how the replay feeds the engines, which data
structures each engine keeps and what it does with every incoming message,
and which settings were fixed for the run. It describes the logic of the
implementation rather than its source code; the clearing step is also given
as pseudocode in \Cref{app:algorithm}.

\section{Orders and Trades}
\label{sec:implorders}

Both engines work on the same two record types. An \emph{order} carries its
order number, the participant that submitted it, its side, its type, its
limit price, its size, its submission time in nanoseconds and the status of
the record it came from. To these the engines add one field of their own,
the \emph{remaining quantity}, which starts at the order's size and is the
only field matching ever changes. A \emph{trade} records the price, the
quantity, the order numbers and participants of the buyer and the seller,
the time, and which engine produced it.

Three conversions are applied when a record of the dataset is read.

\begin{itemize}
  \item \textbf{Prices are integers.} A price is stored as a whole number of
  millionths of a dollar ($10^{-6}$\,USD). All comparisons between prices,
  such as whether a limit price equals the clearing price, are therefore
  exact, with no floating-point rounding.
  \item \textbf{Quantities are whole SOL.} A size is rounded to the nearest
  whole SOL, with halves rounded up. A record whose size rounds to zero, an
  order smaller than 0.5\,SOL, is skipped; this affects \SampleSkippedPct\%
  of all records (\Cref{tab:sample}) and both engines identically.
  \item \textbf{Order types are reduced to two kinds.} Limit, stop-limit and
  take-profit-limit orders become \emph{limit orders} at their limit price.
  Market, stop-market, take-profit-market and vault-close orders become
  \emph{market orders}, which have no limit price.
\end{itemize}

The status of a record decides what the engines do with it
(\Cref{sec:cancellations}): an \texttt{open} record of an ordinary order or
a \texttt{triggered} record of a conditional order becomes a new order, a
record with one of the eight cancel statuses removes the order with the same
number, and every other record is ignored.

\section{The Replay Loop}
\label{sec:implreplay}

The replay reads the hourly files in time order and handles one record at a
time. Each record goes through the same four steps.

\begin{enumerate}
  \item \textbf{Start the clock.} The first record of the run sets the
  origin $t_0$ of the measurement grid \eqref{eq:bucket} and opens the first
  batch, which is due to close at $t_0 + \tau$.
  \item \textbf{Close due batches.} While the record's timestamp is at or
  after the next closing time, the current FBA batch is cleared and the
  closing time moves forward by $\tau$. The FBA therefore has no clock of
  its own: a batch closes when the first record at or after its closing
  time arrives, and after a silence in the data several empty batches are
  closed one after another. Because the check comes before the record is
  handed over, a record stamped exactly at a closing time belongs to the
  next batch, which is the half-open window $[0, \tau)$ of
  \Cref{sec:fbanotation}.
  \item \textbf{Hand the record to both engines.} A new order or a
  cancellation is given first to the FBA, which only stores it, and then to
  the CDA, which matches it at once.
  \item \textbf{Record the state.} The CDA's best bid and ask, their sizes,
  its total depth and any trades are passed to the metrics code after every
  record. The FBA's outcome is passed on at every clear.
\end{enumerate}

Computing time is measured around each CDA submission and around each FBA
clear. These are the only two wall-clock measurements in the run, and they
produce the throughput and clearing-latency controls of
\Cref{sec:metricscontrol}.

\section{The Continuous Engine}
\label{sec:implcda}

\paragraph{Data structures.} Each side of the book is a sorted map from a
price to a queue of the orders resting at that price. The best ask is the
lowest key of the ask map and the best bid the highest key of the bid map,
so both are found without scanning the book. Within a queue, orders are kept
in order of submission time, then order number, which is time priority. An
index from order number to side and price lets a cancellation go straight
to the right queue. Running totals of resting quantity and of the number of
orders on each side are updated with every change, so depth is available at
any moment without adding up the book.

\paragraph{A new order.} For an incoming buy order (a sell is symmetric):
\begin{enumerate}
  \item Look at the first order in the queue of the lowest ask price. If the
  incoming order is a limit order whose price is below that ask, stop.
  \item Otherwise trade $\min(\text{remaining of both})$ at the
  \emph{resting} order's price, stamped with the incoming order's time, and
  reduce both remaining quantities by it.
  \item If the resting order is now fully filled, remove it from its queue
  and from the index, and remove the price level if its queue is empty.
  Repeat from step~1 while the incoming order has quantity left.
  \item If a limit order still has quantity left, it joins the back of the
  queue at its own price. If its timestamp is earlier than that of the last
  order already there, which happens when the data are slightly out of
  order, it is inserted at its time-ordered position instead. Whatever
  remains of a market order is discarded, because a market order never
  rests.
\end{enumerate}

\paragraph{A cancellation.} The index gives the side and price of the order;
the order is removed from that queue, together with the level if it becomes
empty, and the running totals are reduced by its remaining quantity. A
cancellation for an order that is not resting, because it was already
filled or never entered, changes nothing.

The book keeps one structural rule at all times: a price level exists in
the map only while at least one order rests at it. A violation would stop
the run (\Cref{sec:validation}).

\section{The Batch Engine}
\label{sec:implfba}

\paragraph{Data structures.} The FBA keeps a single list of pending orders,
which holds the new orders of the current batch together with everything
rolled over from earlier batches. It also remembers the last price at which
a batch actually traded, $p_{\mathrm{prev}}$, which the tie-break needs.

\paragraph{A new order or a cancellation.} A new order is appended to the
list and nothing else happens until the batch closes. A cancellation
removes the order with that number from the list, including any part of it
that was rolled over; as in the CDA, an unknown order number changes
nothing.

\paragraph{Clearing a batch.} At each close the engine runs the following
steps (\Cref{lst:clearing}).
\begin{enumerate}
  \item \textbf{Candidate prices.} The candidates are the distinct limit
  prices in the list. $D(p)$ and $S(p)$ only change value at these prices,
  so the volume-maximising price is always one of them. If there is no
  candidate, because the list is empty or holds only market orders, nothing
  trades and every order stays in the list.
  \item \textbf{Demand and supply.} The quantities of market buy and sell
  orders are added up once, since they count at every price. Bid limits and
  ask limits are each sorted by price and given a running cumulative sum,
  so $D(p)$ and $S(p)$ at any candidate are read off with one binary search
  instead of a pass over the whole list.
  \item \textbf{Price selection.} The candidates are compared one after
  another under the cascade of \Cref{sec:fbanotation}: larger $V(p)$ first,
  then smaller $|D(p) - S(p)|$, then smaller distance to $p_{\mathrm{prev}}$,
  then the lower price. The winner is $\clr$.
  \item \textbf{Eligible orders.} On each side the orders that accept
  $\clr$ are selected: every market order and every limit order priced at or
  better than $\clr$. They are sorted by priority: market orders first, then
  the better limit price, then the earlier submission time, then the order
  number.
  \item \textbf{Matching.} Two pointers walk down the sorted buy list and
  the sorted sell list together. At each step the two current orders trade
  $\min(\text{remaining of both})$ at $\clr$, and the pointer of whichever
  order is used up moves on. The walk stops when one side runs out, which
  happens after exactly $V^{*} = \min\{D(\clr), S(\clr)\}$ has traded,
  because every eligible buyer accepts every eligible seller at $\clr$. All
  trades of the batch are stamped with the latest submission time among the
  batch's orders.
  \item \textbf{Residual.} Every order with quantity left, whether eligible
  but rationed out or not eligible at all, stays in the list for the next
  batch with its original limit price and submission time, so it keeps its
  time priority. $p_{\mathrm{prev}}$ is updated only if the batch traded.
\end{enumerate}

The cost of a clear grows with the size of the whole list, not only with the
orders that arrived in the last second, because rolled-over orders are
evaluated again in every batch. This is why one FBA clear takes far longer
than one CDA submission (\Cref{sec:res23}).

\section{Fixed Configuration of the Run}
\label{sec:implconfig}

\Cref{tab:config} lists the settings of the run. The batch interval is the
only one that can be changed when the simulator is started; every other
entry is built into the engines, so there is no switch that would select a
different rule.

\begin{table}[htbp]
\centering
\caption[Engine configuration]{Configuration of the run reported in this
thesis.}
\label{tab:config}
\small
\begin{tabularx}{\textwidth}{@{}p{4.2cm} L@{}}
\toprule
\textbf{Setting} & \textbf{Value}\\
\midrule
Engines & CDA and FBA, fed in one pass with the same records\\
Batch interval $\tau$ & 1\,s; also the width of a measurement bucket\\
Batch close & fixed, triggered by event time; no randomisation\\
Candidate prices & the submitted limit prices of the batch\\
Tie-break cascade & volume, then imbalance, then distance to $p_{\mathrm{prev}}$, then lowest price\\
Rationing & price-time priority within the batch\\
Residual & always rolled over; never cancelled by the engine\\
Price unit & $10^{-6}$\,USD (exact integer arithmetic)\\
Quantity unit & 1\,SOL; sizes rounded half-up, below 0.5\,SOL skipped\\
Not implemented & time-in-force, post-only, self-trade prevention, fees and rebates, netting\\
\bottomrule
\end{tabularx}
\end{table}

The missing self-trade prevention is a genuine limitation, not a policy
applied to both engines on purpose (\Cref{sec:experiments}); the effect of
the missing time-in-force and post-only rules is measured in
\Cref{sec:testfidelity}.
